\section{Storytelling}
\label{sec:foundations:storytelling}

Four principles govern the construction of analytical charts intended
to convey a specific finding to a human reader~\cite{knaflic2015}. The
principles are stated as design imperatives rather than as descriptive
observations: each is a rule the designer is expected to apply, not a
property that any well-formed chart automatically possesses.
Table~\ref{tab:foundations:knaflic} summarises the four principles;
the paragraphs that follow elaborate each.

\begin{table}[!htbp]
\centering
\caption{Four principles for analytical chart construction~\cite{knaflic2015}.}
\label{tab:foundations:knaflic}
\footnotesize
\renewcommand{\arraystretch}{1.3}
\begin{tabularx}{\linewidth}{p{2.4cm} X}
\hline
\textbf{Principle} & \textbf{Imperative} \\
\hline
Declutter         & Remove every visual element that does not carry information \\
Focus attention   & Reserve saturated colour and heavy line weight for the element that supports the chart's claim \\
Annotate          & Place a textual interpretation on the chart itself, beside the marks it refers to \\
Structure as narrative & Sequence charts deliberately so they tell a story; the order is part of the design \\
\hline
\end{tabularx}
\end{table}

\textbf{Declutter.} Visual noise competes with signal for the
reader's attention. Gridlines, background fills, axis labels, tick
marks, and legend entries are each candidates for removal and should
be inspected one at a time; absence of an element that does not
impair reading is preferable to its presence.

\textbf{Focus attention.} Default chart elements are rendered in a
single muted tone, with saturated colour and heavy line weight
reserved for the element the analytical claim points at. The principle
inverts the common default in which every category receives a distinct
bright colour and the reader's eye is therefore pulled in many
directions at once. Figure~\ref{fig:foundations:declutter} contrasts a
cluttered chart with its decluttered counterpart on the same data.

\begin{figure}[!htbp]
\centering
\includegraphics[width=0.95\linewidth]{foundations_declutter.png}
\caption{Cluttered and decluttered versions of the same monthly-ridership data.}
\label{fig:foundations:declutter}
\end{figure}

\textbf{Annotate.} An annotation integrated with the marks it
refers to is read together with those marks; a caption beneath the
figure is read as an afterthought. The annotation should state the
finding inline rather than ask the reader to derive it.

\textbf{Structure as narrative.} A sequence of charts that together
tell a story is more effective than a single dense chart that contains
the same information. The order in which charts appear and the
transitions between them is itself part of the design, composed
deliberately rather than left to the order in which the underlying
queries happened to be written. The design space of narrative
visualisation in particular --- author-driven versus reader-driven
ordering, the choice between martini-glass, interactive-slideshow, and
drill-down genres, and the catalogue of visual transitions used to
move from scene to scene --- is mapped out in detail by Segel and
Heer~\cite{segelheer2010}; the four scroll-trigger transition types
adopted in the proposed system's Stories
(Section~\ref{sec:solution:themes:stories}) are a direct application
of their taxonomy.

The four principles are interpretive guides rather than strict rules,
and they occasionally pull against one another: an aggressive
declutter may remove an annotation, and a strong focus-attention
treatment may require the addition of marks that the declutter
principle would otherwise want to remove. The specific resolution of
each tension in the proposed system is recorded in
Section~\ref{sec:solution:interaction}.

Three further structural devices follow from the fourth principle and
are adopted in the present project as concrete implementations of the
narrative recommendation~\cite{knaflic2015}.

The first device is the three-act spine, sometimes called the
\textit{bing-bang-bongo} pattern in the advertising and data-storytelling
literature and adopted by Knaflic~\cite{knaflic2015} for analytical
presentations. Every authored narrative opens with an opening scene that
names the dramatic question (the bing); the middle scenes carry the
chart-by-chart unfolding of the answer (the bang); and the closing scene
recapitulates the finding and names a call to action (the bongo). The three-act spine is required because a
narrative without an opening leaves the reader unsure what question
they are answering, and a narrative without a closing leaves the
reader without a place to act on what they have just read. The
proposed system applies the three-act spine to every story in its
analytical surface (Section~\ref{sec:solution:themes:stories}).

The second device is the action title with supporting subtitle: the
title of a chart states the finding the chart supports, and the
subtitle states the variables and the data window. The action title
carries the analytical point; the subtitle carries the variable labels
and data-window context. The inversion of the conventional title-and-caption order
follows from the storytelling recommendation that the title carry
the conclusion rather than the variable~\cite{knaflic2015}. The device is enforced
uniformly across the catalogue through a shared chart-style helper
described in Section~\ref{sec:solution:interaction:storytelling}.

The third device is the readability gate: a chart is required to
surface its finding to a first-time reader within three to eight
seconds, where eight seconds is the upper bound on the time a reader
typically grants a chart before deciding whether to keep looking. A
chart that fails the gate is re-designed before being added to the
catalogue. The gate is operationalised through a glance protocol: a
reader unfamiliar with the chart looks at it for five seconds, looks
away, and articulates the finding; a failure to articulate signals a
re-design.

The three devices share a common goal: the chart and the narrative
should each lead the reader to the same finding without verbal
mediation. The proposed system implements them as enforceable
constraints on chart construction and on narrative authoring, not as
recommendations to be left to the discretion of the designer.
